\section{Rotation and return to the original manifold}
\label{sec:rotation}

We now complete the proof of Theorem~\ref{thm:main}.  The marking in
Proposition~\ref{prop:family} will be used throughout: all second
cohomology groups are identified with the lattice of the differentiable
manifold underlying the given $X$, and $q$ denotes both the quadratic
form and its associated symmetric bilinear form.

\subsection{The primitive class and the chosen fiber}

Recall the primitive class $e$ and line bundle $E$ constructed at the
start of Section~\ref{sec:family}. Write $E_X=E$, so that
$L\simeq E_X^{\otimes m_0}$. The class $e$ is nonzero, nef, and
$q(e)=0$; its sign was fixed so that, for every K\"ahler class
$\kappa$ on $X$,
\begin{equation}
                    q(e,\kappa)>0.
\label{eq:rotation-original-sign}
\end{equation}

Fix $\lambda\in\R\setminus\Q$ and use the family of
Proposition~\ref{prop:family}.  Its rational classes $e,f,v,h$ satisfy
\[
 q(e)=q(f)=0,\quad q(e,f)=1,\quad
 q(e,v)=q(f,v)=0,\quad a:=q(v)>0,
\]
and $h\perp P:=\langle e,f,v\rangle_{\Q}$.  The normalized period
and the local monodromy $M$ on $H^2$ are
\begin{equation}
 p(\tau)=f+\tau v-\frac a2\tau^2e,
       \qquad\tau=\lambda+iy,\quad y>0,
\label{eq:rotation-period}
\end{equation}
\begin{equation}
 Me=e,\qquad Mv=v-aw e,\qquad
 Mf=f+w v-\frac{aw^2}{2}e,\quad w\ne0.
\label{eq:rotation-monodromy}
\end{equation}
Consequently,
\begin{equation}
                    (M-1)P=\langle e,v\rangle_{\Q}.
\label{eq:rotation-monodromy-image}
\end{equation}
Theorem~\ref{thm:torus}, applied along the prescribed vertical ray,
provides a smooth fiber $Y=Y_t$, its polarized Ricci-flat metric $g$,
and an embedded special Lagrangian torus $T\subset Y$.  Here
$\dim_{\C}Y=2n$, $\dim_{\R}T=2n$, and $T$ is isotopic to a
coordinate torus in a semistable chart
$z_0z_1\cdots z_{2n}=t$.  A constant rescaling of $g$, as used in
the torus construction, does not affect these properties.

Write the present complex structure as $I$.  Choose a hyperk\"ahler
triple $(I,J,K)$ for $g$, with $IJ=K$, so that
\begin{equation}
       \omega_J+i\omega_K=c\sigma_t,\qquad c>0,
       \qquad [\omega_I]\in\R_{>0}h.
\label{eq:rotation-normalization}
\end{equation}
The positive scalar normalizes the length of $\sigma_t$; its given
phase fixes the ordered pair $(J,K)$.  Since $[\sigma_t]=p(\tau)$,
\begin{equation}
 [\omega_K]=cy(v-a\lambda e),\qquad
 q(e,[\omega_J])=c>0.
\label{eq:rotation-k-class}
\end{equation}
The second equality follows from $q(e,p(\tau))=1$ and will retain
the required sign of $e$ after rotation.

\subsection{Restriction, calibration, and complex rotation}

The following criterion isolates the rotation argument. Cohomological
vanishing is enough because the special Lagrangian condition turns
volume into a period.

\begin{lemma}[A cohomological rotation criterion]
\label{lem:rotation-criterion}
Let $(Y,g,I,J,K)$ be a hyperk\"ahler manifold of real dimension $4n$,
where $n\ge1$ and $IJ=K$, and let $i:T\hookrightarrow Y$ be a compact,
connected, oriented embedded
special Lagrangian submanifold for the complex structure $I$.
If $i^*[\omega_K]=0$ in $H^2(T,\R)$, then $T$ is a complex
$n$-dimensional submanifold for $J$ and is Lagrangian for
$\omega_K+i\omega_I$. If $T$ is diffeomorphic to a real $2n$-torus,
then it is a complex torus for $J$.
\end{lemma}
\begin{proof}
The normalized parallel holomorphic volume form for $(Y,I,g)$ is
\[
                 \Omega_I=\frac{(\omega_J+i\omega_K)^n}{n!}.
\]
To check the normalization, in a symplectic unitary frame one has
$\omega_J+i\omega_K=\sum_{j=1}^n dz_{2j-1}\wedge dz_{2j}$;
its $n$th power divided by $n!$ is the unit holomorphic volume
form.  It has absolute value one on an orthonormal frame of any
$I$-Lagrangian real $2n$-plane. The special Lagrangian calibration identity
\cite[III.1, Theorem~1.10 and Corollary~1.11]{HarveyLawson}
therefore gives
\[
                    \vol_g(T)=\left|\int_T\Omega_I\right|.
\]
All terms containing $\omega_K$ in this integral vanish:
the hypothesis makes $i^*\omega_K$ exact,
while all the factors in the expansion are closed.  Stokes' theorem
then yields
\begin{equation}
 \vol_g(T)
 =\left|\int_T\frac{\omega_J^n}{n!}\right|
 \le\int_T\left|\frac{i^*\omega_J^n}{n!}\right|
 \le\vol_g(T).
\label{eq:rotation-wirtinger}
\end{equation}
The last inequality is the pointwise Wirtinger inequality
\cite[II.6, Theorem~6.11]{HarveyLawson}.
Its continuous, nonnegative deficit has integral zero, so equality
holds at every point of $T$.  Equality planes are exactly the
$J$-invariant real $2n$-planes.  Thus $T$ is a complex
$n$-dimensional submanifold of
\[
                              Y'=(Y,J).
\]
Its orientation may now be chosen to be the complex orientation;
the absolute period in \eqref{eq:rotation-wirtinger} made the
initial choice immaterial.

For $J$ the holomorphic symplectic form is
$\sigma_J=\omega_K+i\omega_I$.  We already have
$\omega_I|_T=0$ pointwise.  Since $JT_xT=T_xT$, for
$u,v\in T_xT$ the quaternion identities give
\[
              \omega_K(u,v)=g(IJu,v)=\omega_I(Ju,v)=0.
\]
Hence $\sigma_J|_T=0$: the rotated submanifold is holomorphically
Lagrangian.

Suppose now that $T$ is diffeomorphic to a real $2n$-torus. The
restriction of $\omega_J$ makes $T$ K\"ahler.
Consequently $b_1(T)=2n$ and $h^{1,0}(T)=n$.  Its Albanese map
\[
                 \alpha:T\longrightarrow\operatorname{Alb}(T)
\]
has a complex $n$-torus as target and induces an isomorphism on
integral first cohomology.  Both integral cohomology rings are
exterior algebras on first cohomology, so $\alpha^*$ is an
isomorphism in every degree.  In particular, $\alpha$ has
nonzero degree; holomorphicity makes that degree positive, hence
equal to one, and $\alpha$ is surjective.

Choose a K\"ahler class $\xi$ on $T$.  Surjectivity on $H^2(-,\R)$
writes $\xi=\alpha^*\beta$ for some real class $\beta$ on
$\operatorname{Alb}(T)$.  If a fiber contained a positive-dimensional
irreducible compact analytic subvariety $V$ of dimension $k$, then
\[
                        \int_V\xi^k=0,
\]
since $\alpha|_V$ is constant.  This contradicts K\"ahler
positivity.  Thus $\alpha$ is proper with finite fibers, hence
finite.  A finite degree-one holomorphic map to a normal complex
space is an isomorphism.  This proves that $T$ is a holomorphic
Lagrangian complex torus in $Y'$.
\end{proof}

We verify the criterion for the torus constructed above, using its
coordinate isotopy class. This is where the monodromy information is
needed.

For a coordinate torus $T_0=(\R/2\pi\Z)^{2n}$ choose its fixed
radii $r_1,\ldots,r_{2n}$.  Around the circle
$t(\theta)=t\exp(i\theta)$, the formula
\[
 z_j=r_j\exp(iy_j)\quad(1\le j\le2n),\qquad
 z_0=\frac{t\exp(i\theta)}{z_1\cdots z_{2n}}
\]
defines a family of embeddings of this same parameterized $T_0$.
All absolute values remain fixed, so the family stays inside the
chart, and the embeddings at $\theta=0$ and $2\pi$ agree as
parameterized maps.  The induced restriction maps form a morphism
from the fiber cohomology local system to the constant local system
$H^k(T_0,\R)$.  Therefore, in every degree $k$,
\begin{equation}
 i_0^*M_k=i_0^*:H^k(Y,\R)\longrightarrow H^k(T_0,\R),
\label{eq:rotation-restriction-monodromy}
\end{equation}
where $M_k$ is monodromy in degree $k$.  This statement controls
the restriction of classes in degree two, as well as the torus's
top-degree period.  Isotopy identifies the restriction maps for
$T_0$ and $T$.  Thus, for $i:T\hookrightarrow Y$,
$i^*(M-1)=0$ on $H^2(Y,\R)$.  Equations
\eqref{eq:rotation-monodromy-image} and
\eqref{eq:rotation-k-class} give
\begin{equation}
                         i^*[\omega_K]=0.
\label{eq:rotation-k-vanishing}
\end{equation}
Replacing monodromy by its inverse or a positive power preserves
the image in \eqref{eq:rotation-monodromy-image}, so the conclusion
is independent of this convention and of the finite base changes.

Lemma~\ref{lem:rotation-criterion} now makes $T$ a holomorphic
Lagrangian complex torus on the rotated manifold $Y'=(Y,J)$.

\subsection{The rational isotropic ray after rotation}

The period plane of $Y'$ is
\begin{equation}
       \langle[\omega_K],[\omega_I]\rangle_{\R}
       =\langle v-a\lambda e,h\rangle_{\R}.
\label{eq:rotation-j-period}
\end{equation}
It is perpendicular to $e$, so $e$ is an integral $(1,1)$ class
on $Y'$; see Figure~\ref{fig:rotation}.

\begin{figure}[htbp]
\centering
\begin{tikzpicture}[scale=.92,>=Stealth,font=\small]
  \coordinate (O) at (0,0);
  \coordinate (I) at (3.5,-.65);
  \coordinate (J) at (-3,-.9);
  \coordinate (K) at (0,2.65);
  \fill[blue!9] (O)--(-2.35,-.705)--(-2.35,1.35)--(0,2.055)--cycle;
  \fill[orange!14] (O)--(2.65,-.492)--(2.65,1.563)--(0,2.055)--cycle;
  \draw[->,semithick] (O)--(I) node[right,align=left] {$[\omega_I]$\\$\in\R_{>0}h$};
  \draw[->,semithick] (O)--(J) node[left] {$[\omega_J]$};
  \draw[->,semithick] (O)--(K) node[above] {$[\omega_K]\in\R_{>0}(v-a\lambda e)$};
  \node[align=center] at (-1.27,.78) {period plane\\for $I$};
  \node[align=center] at (1.42,.90) {period plane\\for $J$};
  \fill (O) circle (1.2pt);
\end{tikzpicture}
\caption{The period planes before and after rotation, inside the
positive three-space spanned by the hyperk\"ahler forms.
The class $e$ is perpendicular to $[\omega_I]$ and $[\omega_K]$,
so it becomes a $(1,1)$ class for $J$; its pairing with
$[\omega_J]$ fixes the positive sign.}
\label{fig:rotation}
\end{figure}


Let
$x\in H^{1,1}(Y')\cap H^2(Y',\Q)$.  Orthogonality to the period
plane says
\[
                q(x,h)=0,\qquad q(x,v)=a\lambda q(x,e).
\]
The three numbers $q(x,v)$, $a$, and $q(x,e)$ are rational,
whereas $a\ne0$ and $\lambda\notin\Q$.  Therefore
\begin{equation}
             q(x,e)=q(x,v)=0.
\label{eq:rotation-rational-orthogonality}
\end{equation}

The restriction of $q$ to $H^{1,1}(Y',\R)$ has signature
$(1,b_2-3)$.  In a Lorentzian space, the orthogonal complement of
a nonzero null vector is negative semidefinite with radical its
own line.  Applying this to $e$ and
\eqref{eq:rotation-rational-orthogonality} proves
\begin{equation}
 \begin{gathered}
 q(x)\le0\quad\text{for }x\in H^{1,1}(Y')\cap H^2(Y',\Q),\\
 q(x)=0\quad\Longrightarrow\quad x\in\Q e.
 \end{gathered}
\label{eq:rotation-unique-ray}
\end{equation}
In particular, $Y'$ is nonprojective: a projective hyperk\"ahler
manifold has an integral ample class of positive BBF square.
This is the necessary direction of the projectivity criterion
\cite{Huybrechts,HuybrechtsErratum}.  The rational Picard space
may contain negative-square classes; its isotropic classes are
exactly those on $\Q e$.  Equation~\eqref{eq:rotation-k-class}
places $e$ in the closure of the positive cone containing the
K\"ahler class $[\omega_J]$.

\subsection{A semiample point and deformation back to \texorpdfstring{$X$}{X}}

We state explicitly the two external results used in this step.

\begin{theorem}[Greb--Lehn--Rollenske {\cite[Theorem~4.1]{GLR}}]
\label{thm:nonprojective-torus}
If a nonprojective compact irreducible hyperk\"ahler manifold $Z$ of
complex dimension $2n$ contains a holomorphic Lagrangian complex
torus $S$, then its algebraic dimension is $n$ and it has a
holomorphic algebraic reduction $Z\to B$ which is a Lagrangian
fibration with $S$ as a fiber.  In particular, this fibration has
normal projective base and provides a nontrivial semiample
pullback bundle on $Z$.
\end{theorem}
The projectivity assertion follows from the line-bundle construction
of the algebraic reduction in \cite[Section~4.2]{GLR} and Stein
factorization.

\paragraph{The fixed-class deformation theorem.}
Let $N$ be a compact simply connected differentiable manifold
of hyperk\"ahler type with maximal holonomy, let
$\mathcal T^\circ(N)$ be one connected component of its
Teichm\"uller space, and fix a nonzero primitive integral class
$e$ with $q(e)=0$.  Write $\mathcal C_I$ for the component of
the positive cone containing the K\"ahler cone of $(N,I)$, and set
\[
 \mathcal T^\circ_e
 =\{[I]\in\mathcal T^\circ(N):
       e\in H^{1,1}(N,I),\ e\in\overline{\mathcal C_I}\}.
\]
This is the positive-sign locus of
\cite[Definition~2.6]{SV}.  By \cite[Proposition~2.7]{SV} it is
connected.  Let $E_I$ denote the unique holomorphic bundle with
class $e$.  Define $\mathcal T^\circ_{e,\mathrm{nef}}$ and
$\mathcal T^\circ_{e,\mathrm{sa}}$ by requiring $E_I$ to be nef
and semiample, respectively.

\begin{theorem}[Soldatenkov--Verbitsky {\cite[Theorem~3.7(v)]{SV}}]
\label{thm:fixed-class}
With this notation,
\begin{equation}
 \mathcal T^\circ_{e,\mathrm{sa}}\ne\varnothing
 \quad\Longrightarrow\quad
 \mathcal T^\circ_{e,\mathrm{nef}}
      =\mathcal T^\circ_{e,\mathrm{sa}}.
\label{eq:rotation-sv}
\end{equation}
\end{theorem}
Its family formulation, \cite[Theorem~3.8]{SV}, says that a nef
isotropic bundle is semiample when its pair deforms, in a smooth
connected family carrying a line bundle, to a semiample pair.
We use \eqref{eq:rotation-sv} directly.

Apply Theorem~\ref{thm:nonprojective-torus} to $(Y',T)$.  Choose an ample bundle $H$
on its fibration base and let $A$ be its pullback to $Y'$.
The bundle $A$ is semiample and nef, and its first Chern class
$b$ is nonzero: a semipositive representative of a generated
power has positive trace on an open set where the map has
positive rank.  Since the base has dimension $n$,
$\int_{Y'}b^{2n}=0$.  Fujiki's relation \cite{Fujiki} gives
$q(b)=0$.  Equation~\eqref{eq:rotation-unique-ray} and the
primitivity of $e$ now give
\[
                             b=k e,\qquad k\in\Z.
\]
Both $b$ and $e$ pair positively with $[\omega_J]$: for $b$ this
is the nonzero nef null-vector property, and for $e$ it is
\eqref{eq:rotation-k-class}.  Hence $k>0$.

Let $E_{Y'}$ be the unique bundle on $Y'$ with first Chern class
$e$.  Injectivity of the first Chern class map gives
\[
                          A\simeq E_{Y'}^{\otimes k}.
\]
A globally generated positive power of $A$ is therefore a
globally generated positive power of $E_{Y'}$.  Thus the
\emph{specified primitive class} $e$ is semiample on $Y'$.

The marking furnished by Proposition~\ref{prop:family} identifies
$Y$ with a point of the same ambient component
$\mathcal T^\circ(N)$ as the original $X$.  The hyperk\"ahler
twistor sphere connects $I$ to $J$ on that same marked
differentiable manifold, so $Y'$ also belongs to this component.
On $X$, $e$ is of type $(1,1)$ and has the positive sign by
\eqref{eq:rotation-original-sign}.  On $Y'$ the corresponding
statements follow from \eqref{eq:rotation-j-period} and
\eqref{eq:rotation-k-class}.  Consequently both endpoints lie
in the same positive fixed-class locus $\mathcal T^\circ_e$.
The semiample point $Y'$ makes its semiample locus nonempty;
the original point $X$ belongs to its nef locus by hypothesis.
Equation~\eqref{eq:rotation-sv} now proves that $E_X$ is
semiample on $X$.

Choose $s>0$ such that $E_X^{\otimes s}$ is globally generated.
Then $L^{\otimes s}\simeq(E_X^{\otimes s})^{\otimes m_0}$ is
globally generated as well.  In particular, on the original
manifold the evaluation morphism
\[
 H^0(X,L^{\otimes s})\otimes_{\C}\mathcal O_X
                    \longrightarrow L^{\otimes s}
\]
is surjective everywhere.  This proves the semiampleness assertion
of Theorem~\ref{thm:main}.

\subsection{The fibration defined by a generated power}

We derive the fibration description from semiampleness, using the
Fujiki argument for Lagrangian fibers of Matsushita
\cite{Matsushita1999,MatsushitaAddendum}.  Let $m>0$ make $L^{\otimes m}$ globally
generated, let
\[
 g:X\longrightarrow Z\subset\mathbb P H^0(X,L^{\otimes m})^*
\]
be its morphism to the image, and take its Stein factorization
$g=\nu\circ f$:
\[
                   X\xrightarrow{\ f\ }B\xrightarrow{\ \nu\ }Z.
\]
The space $B$ is normal, $f$ has connected fibers, and $\nu$ is
finite.  Since $Z$ is projective, $B$ is projective and
$H:=\nu^*\mathcal O_Z(1)$ is ample.  By construction
\begin{equation}
                         L^{\otimes m}\simeq f^*H.
\label{eq:rotation-ample-pullback}
\end{equation}

Put $\ell=c_1(L)$ and choose a K\"ahler class $\kappa$ on $X$.
The positive sign gives $q(\ell,\kappa)>0$.  Fujiki's formula
and $q(\ell)=0$ give the polynomial identity
\begin{equation}
 \int_X(\kappa+t\ell)^{2n}
       =c_X\bigl(q(\kappa)+2tq(\kappa,\ell)\bigr)^n,
       \qquad c_X>0.
\label{eq:rotation-numerical-dimension}
\end{equation}
Comparison of coefficients shows
\[
 \int_X\ell^k\kappa^{2n-k}
 \begin{cases}
 >0,&0\le k\le n,\\
 =0,&n<k\le2n.
 \end{cases}
\]
On the other hand, a pullback of the Fubini--Study form by $g$
represents $m\ell$, is semipositive, and has generic rank
$\dim Z=\dim B$.  Its $k$th power wedged with a K\"ahler form
has positive integral precisely for $k\le\dim B$.  Thus
\begin{equation}
                              \dim_{\C}B=n.
\label{eq:rotation-base-dimension}
\end{equation}

Let $F$ be a general smooth fiber.  Its Poincar\'e dual class
is a positive multiple of $\ell^n$.  Indeed, choose a very ample
power $H^{\otimes a_0}$ and intersect $n$ general members of
its linear system on $B$.  Their intersection is a nonempty
finite set of transverse points in the smooth regular-value
locus of $f$.  Pulling back gives the union of the corresponding
smooth fibers and represents $(a_0m)^n\ell^n$.  These fibers
have the same homology class, by smooth proper transport along
the connected regular-value locus.  This proves the claim.

For $n\ge2$, polarization of Fujiki's relation gives
\begin{equation}
              \int_X\ell^n\kappa^{n-2}\sigma\overline\sigma=0.
\label{eq:rotation-fiber-vanishing}
\end{equation}
To see the vanishing directly, each summand in the polarized
formula is a product of pairings of the $2n$ degree-two factors.
Every copy of $\ell$ is orthogonal to $\ell$, $\sigma$, and
$\overline\sigma$, so it would have to pair with a copy of
$\kappa$.  There are $n$ copies of $\ell$ and only $n-2$ copies
of $\kappa$, and therefore every summand is zero.
Using the fiber class, \eqref{eq:rotation-fiber-vanishing} gives
\[
              \int_F\sigma|_F\wedge\overline{\sigma|_F}
                         \wedge\kappa|_F^{\,n-2}=0.
\]
In unitary coordinates this integrand is a fixed positive
constant times the squared pointwise norm of $\sigma|_F$
times the K\"ahler volume form.  Hence $\sigma|_F=0$.
For $n=1$, the restriction of a holomorphic two-form to the
curve $F$ vanishes automatically.  The fibers in question
have complex dimension $n$ by
\eqref{eq:rotation-base-dimension}, so they are Lagrangian.

The same assertion holds for every irreducible component of every
fiber by Matsushita's equidimensionality theorem
\cite[Theorem~1]{MatsushitaEquidimensional}. Its hypotheses are
exactly those now established: $X$ is K\"ahler, $f$ is proper and
surjective with normal base, and a general smooth fiber is
Lagrangian. The holomorphic symplectic form is closed, as every
holomorphic form on a compact K\"ahler manifold is. Thus every
fiber component has dimension $n$ and is isotropic on its smooth
locus. This also gives the Lagrangian fibration used in
Section~\ref{sec:finiteness}.

Conversely, an ample bundle has a globally generated positive
power, so \eqref{eq:rotation-ample-pullback} implies semiampleness
of $L$. This proves the equivalence between semiampleness and the
fibration description in Theorem~\ref{thm:main}.

\subsection{Projective-space bases}
\label{subsec:projective-base}

We finish with the additional conclusion for projective $X$ in
Corollary~\ref{cor:projective-base}.

\begin{proof}
Choose a very ample line bundle $H$ on $X$, independently of $L$, with
embedding $i:X\hookrightarrow\mathbb P^N$; the closed graph embedding
$(f,i):X\hookrightarrow B\times\mathbb P^N$ exhibits $H$ as relatively
very ample for $f$, so $f$ is projective. The remaining hypotheses are
supplied by Theorem~\ref{thm:main}; hence the companion's
Theorem~1.1~\cite{ProjectiveSpaceBases2026} identifies $B$ with
$\mathbb P^n_{\C}$, and the following no-multiple-fiber paragraph there
gives the divisor assertion.
\end{proof}
